\section{Literature Review}
\label{sec:literature}
Research on data-centre flexibility spans workload scheduling, electrical storage, cooling control and coordinated facility operation. These strands establish the potential of individual resources and the importance of managing their interactions. The remaining assessment question is how to express the resulting facility response in terms of a requested change at the grid connection, its duration and its consequences for subsequent operation.

IT workload is a principal source of controllable demand. Cao et al.~\cite{Cao2022ApEn} investigate the temporal shifting capability of periodic batch workloads using a data-driven flexibility assessment. The approach makes workload characteristics central to the available response rather than treating the entire IT load as freely dispatchable. Misaghian et al.~\cite{Misaghian2023TIA} assess carbon-aware flexibility measures using machine learning, while Radovanovi\'c et al.~\cite{radovanovic2022carbon} describe carbon-aware computing that shifts eligible work in response to forecast grid conditions. These studies demonstrate the value of distinguishing work arrival from execution. They also show why flexibility depends on which tasks may move and which obligations must be met, not simply on the level of IT demand.

UPS batteries provide a complementary form of flexibility. Their power exchange can change without rescheduling the computational task being supplied. Alaper\"a et al.~\cite{alapera2018data} examine data centres as a source of dynamic flexibility, and subsequent pilot work investigates fast frequency response from a UPS system~\cite{alapera2019fast}. Such applications motivate the use of existing electrical storage for services beyond backup. They do not remove the need to reserve energy for the UPS's primary role or to consider conversion losses and battery ageing. A conventional storage formulation is appropriate when the research question concerns the battery's interaction with the rest of the facility rather than a new representation of battery physics.

Cooling flexibility is linked to heat production and to the thermal state of the facility. Cupelli et al.~\cite{cupelli2018data} develop a data-centre control framework for demand response using models of the coupled subsystems. Fu et al.~\cite{Fu2020ApEn} investigate coordinated IT and cooling control for frequency regulation. Thermal storage extends these opportunities by allowing cooling to be produced before it is needed, although its usefulness remains bounded by storage energy, charge and discharge rates, and the available chiller capacity. Work on data-centre heat recovery and short-term thermal storage also illustrates the importance of distinguishing the timing of thermal production and demand~\cite{du2025new}. The present study concerns cooling within the facility, rather than a district-heating export system.

Integrated operation is thus an established research direction. Cioara et al.~\cite{cioara2018optimized} consider flexibility management for data-centre participation in demand-response programmes, and Xu et al.~\cite{Xu2023SPIE} assess demand-response potential across multiple flexible resources. Cupelli et al.~\cite{cupelli2018data} are particularly relevant because the coupling between IT demand, electrical storage and thermal dynamics is already central to their control framework. The contribution sought here is therefore not the first combination of those resources. It is an explicit assessment of the power deviations and durations that their coordinated operation can sustain from a specified operating state, together with an account of how the components deliver the response and recover afterwards.

A scheduled operating profile and a flexibility envelope answer different questions. Cost or carbon optimisation identifies a preferred trajectory under its objective. A flexibility assessment asks which departures from that trajectory remain feasible. In a storage- and workload-constrained facility, feasibility during an event is insufficient if the response leaves depleted stores, excessive temperatures or additional unfinished work beyond the assessment window. Recovery conditions make these consequences part of the service definition. They also expose why the same requested grid deviation can require different internal actions at different times.

This study combines that duration-oriented assessment with annual economic evaluation. A single daily electricity-price profile can illustrate scheduling behaviour, but cannot establish the annual saving associated with repeated operation. Linked daily optimisation supplies that broader economic evidence while preserving the operating state used by the event assessment. The analysis retains the distinction between electricity-cost savings and grid-service revenues: the former are evaluated over the year, whereas the latter would require market-specific prices, availability requirements and delivery rules beyond the physical capability assessment conducted here.
